\introChapter{ตรรกศาสตร์ (Logic)}

\begin{description}[leftmargin=2cm, style=nextline, itemsep=0.1em, parsep=0pt]
    \item[รหัสวิชา] 4091606
    \item[ชื่อวิชา] คณิตศาสตร์สำหรับคอมพิวเตอร์
    \item[หน่วยกิต] 3 (3-0-6)
    \item[บทที่] 1
    \item[ชื่อบท] ตรรกศาสตร์ (Logic)
    \item[จำนวนชั่วโมง] 6 ชั่วโมง
\end{description}

\noindent \textbf{\large วัตถุประสงค์การเรียนรู้ประจำบท}\\[0.2em]
\noindent เมื่อนักศึกษาเรียนบทนี้จบแล้ว นักศึกษาจะสามารถ
\begin{enumerate}[topsep=0.2em, itemsep=0.15em, leftmargin=*]
    \item อธิบายความหมายและสมบัติของประพจน์ และค่าความจริงได้
    \item ใช้ตัวเชื่อมทางตรรกศาสตร์พื้นฐาน 5 ตัว ได้แก่ $\neg$, $\land$, $\lor$, $\rightarrow$, $\leftrightarrow$ และอธิบายตัวเชื่อมเพิ่มเติม เช่น $\oplus$ ได้อย่างถูกต้อง
    \item สร้างและวิเคราะห์ตารางค่าความจริงของประพจน์เชิงประกอบได้
    \item ตรวจสอบสมมูลทางตรรกศาสตร์ และใช้กฎพีชคณิตของประพจน์ได้
    \item อธิบายและประยุกต์ใช้การอนุมาน เช่น Modus Ponens, Modus Tollens ได้
    \item เชื่อมโยงหลักการตรรกศาสตร์กับการเขียนโปรแกรมและวงจรดิจิทัลได้
\end{enumerate}

\noindent \textbf{\large หัวข้อที่สอน}\\[0.2em]
\begin{enumerate}[topsep=0.2em, itemsep=0.15em, leftmargin=*]
    \item ประพจน์
    \item ตัวเชื่อมทางตรรกศาสตร์
    \item ตารางค่าความจริง
    \item สัจนิรันดร์และข้อขัดแย้ง
    \item สมมูลทางตรรกศาสตร์และกฎพีชคณิตของประพจน์
    \item ประพจน์เลือกเฉพาะ (XOR)
    \item ประพจน์เปิดและตัวบ่งปริมาณ
    \item กฎการอนุมาน
\end{enumerate}

\noindent \textbf{\large สื่อการสอน}
\begin{enumerate}[topsep=0.2em, itemsep=0.15em, leftmargin=*]
    \item สไลด์ประกอบการสอน
    \item แบบฝึกหัดท้ายบท
\end{enumerate}
